The child had begun to move in earnest, and Cersei Lannister had discovered, somewhat to her own irritation, that there was no controlling it.

She had controlled a great many things in her life, by force or guile or simple refusal to accept that a thing could not be bent to her will. She had controlled kings and councils and the terror of an entire city. She had not been able to control the tremble in her own hands the night the Sept of Baelor burned, though she had hidden it well enough that no one living knew of it, and she could not control this — the small, insistent turning beneath her ribs that woke her in the cold hours and would not be reasoned with, a Lannister already, she thought, in its refusal to be told what to do.

She was too old for this. The maesters did not say so to her face, because the maesters who said such things to Cersei Lannister’s face did not remain maesters for very long, but she saw it in the careful way Qyburn watched her rise from a chair now, in the extra beat before he answered when she asked him whether all was as it should be. She was too old, and the winter was too hard, and the child came anyway, indifferent to both, the way her children had always come — Joffrey, Myrcella, Tommen, each of them arriving into a world that had eventually killed them, each of them hers in a way nothing else had ever been.

She would not lose this one. She had decided that with the particular flat certainty she reserved for the few things in her life she refused to leave to the gods. Whatever else the war took — the throne, the city, Jaime, her own life if it came to that — it would not take this child. She had buried three. The gods, or fate, or the spiteful machinery of prophecy that had haunted her since a tent in the woods thirty years past, had taken three. This one she would keep, if she had to burn the whole of Westeros to do it.

\sceneBreak

Qyburn brought her the northern rider’s report in the morning, and Cersei read it twice, slowly, because it was the first genuinely good news the war had brought her in longer than she cared to count.

“The army that broke against the dead at Winterfell,” Qyburn summarized, in his mild, papery voice, “is considerably less than the army that marched north. Half the Dothraki dead, or near it. The Unsullied bloodied badly. The Northern levies starving, and deserting more each night. Two dragons lost — one to the enemy, one grounded and killed, both now serving the dead rather than the girl who hatched them. She has one dragon left, Your Grace. One. And a column of hungry, frightened people limping south toward us with the dead at their backs, in no condition to besiege a fed city behind strong walls.”

“One dragon.” Cersei turned the words over the way she’d once turned over the news of Daenerys Targaryen’s landing — with the wariness of a woman who had learned, expensively, not to celebrate a rival’s weakness prematurely. “One dragon burned the Tarly host to ash and took the Reach in an afternoon, Lord Qyburn. I’d not grow careless about one dragon.”

“Nor would I counsel it, Your Grace. Only that one dragon is a considerably more manageable problem than three, and a dragon whose roosting-place we now know is more manageable still.” Qyburn’s small smile held the particular satisfaction he reserved for problems he believed he had solved. “The scorpions Wisdom Hallyne’s designs left us have been improved upon. I’ve had them mounted along the walls and on the largest ships in Lord Greyjoy’s fleet these past weeks. A dragon that must land to rest is a dragon that must, at some point, be somewhere a bolt can reach it.”

Cersei said nothing for a moment, one hand resting unconsciously where the child turned beneath her ribs. She thought of Daenerys Targaryen — a girl she had never met, and hated with the particular purity she reserved for people who wanted what was hers — limping south with one dragon and a starving army, believing, no doubt, that she was marching toward a reckoning she would win. Cersei had believed such things herself, at various points, and had learned each time that the reckoning rarely arrived shaped the way the marcher expected.

“The dead,” she said. “Behind them. How close?”

“Close enough that the column cannot stop to rest properly. Not so close that they’ll arrive together, by the rider’s estimate. The Dragon Queen’s army will reach us first — days ahead of what pursues them, perhaps a fortnight.” Qyburn paused. “Which raises a question, Your Grace, that I confess I have not been able to answer to my own satisfaction. When they arrive — starving, and desperate, and with the dead at their backs — what is it you intend to do with them?”

\sceneBreak

It was the right question, and Cersei had been turning it over herself for days, and had not yet said the answer aloud to anyone, because saying it aloud made it real in a way she was not quite ready for.

The obvious course was alliance. Even she could see it — a strategist warmer than herself would have opened the gates, joined the two armies, faced the dead as one host behind the strongest walls in Westeros. Tyrion would counsel it, if Tyrion were here. Even Jaime, gods help him, would counsel it, which was half of why he’d gone north in the first place, the great fool, to stand for something in a war she’d already decided was not hers to bleed for.

But Cersei had done the arithmetic that the warmer strategists never quite let themselves finish. If she opened the gates, and the two armies became one, and they somehow won — then Daenerys Targaryen stood victorious in Cersei’s city, with the North and the Reach and the Unsullied and a dragon behind her, and Cersei Lannister became, at best, a defeated rival kept alive out of a conqueror’s mercy, which was to say not kept alive very long at all. She had seen what happened to the losers of wars. She had made it happen, more than once. She did not intend to be on the receiving end of it.

And there was the other arithmetic, the one she liked considerably better. Two enemies were marching toward her, and one of them was going to arrive first, exhausted and starving and desperate, and the second was going to arrive behind it — and if Cersei simply held her walls, and let the dead and the Dragon Queen’s army break against each other outside her gates, then she might, with a little patience and a great deal of wildfire, watch both her enemies destroy each other while she waited, fed and warm and safe, behind the strongest walls in the world.

“I intend to let them come,” she said at last. “And I intend to keep my gates shut when they do. Let the dead and the dragon’s rabble decide it between themselves outside my walls. Whichever wins, they’ll be too spent afterward to trouble a city that spent the battle resting behind stone. That’s how you win a war you were never strong enough to win outright, Lord Qyburn. You let your enemies do your fighting for you, and you keep the last army in Westeros for yourself.”

Qyburn inclined his head, unsurprised, perhaps even approving. “A cold strategy, Your Grace.”

“The warm ones get you killed. I’ve buried three children learning that.” Cersei rose, slowly, one hand at her belly and the other on the chair for balance she resented needing, and crossed to the window that looked out over her city — her city, still, whatever the girl with one dragon believed she was marching toward. “See to the wildfire. All of it. Every cache Hallyne’s boy can pour. If it comes to the walls being tested, I’ll not have this city taken by dead men or dragon-queens or anyone else. I’d sooner it burned. I’d sooner all of it burned than see it kneel to her.”

\sceneBreak

She did not entirely mean the last of it. That was the thing she would not have admitted to Qyburn, or to anyone — that the words had come out of her more easily than they should have, the old Targaryen madness of \textit{burn them all} surfacing in a Lannister mouth, and that some small, cold part of her had meant it entirely, and that the meaning had frightened even her, a little, in the moment it left her.

She had stood in this city as it nearly burned once, when Aerys’s caches went up beneath the Sept and took her enemies and half her own reason with them. She had told herself, then, that she had done what was necessary. She told herself the same now, watching her city go about its business below, unaware how easily its queen had just spoken of burning it to keep it.

The child turned beneath her ribs, insistent, alive, hers. Cersei rested her hand against it and thought of Maggy the Frog, thirty years dead, in a tent that smelled of blood and goat. \textit{Gold shall be their crowns, and gold their shrouds.} Three children buried, all of them gold-haired, all of them crowned in one fashion or another before the shrouds came. And the last of the prophecy, the part she had spent thirty years assigning to Tyrion because Tyrion was easier to fear than the alternative — \textit{the valonqar shall wrap his hands about your pale white throat.}

She was too clever, in the end, not to know which little brother the prophecy had likely always meant. She simply refused, as she had always refused, to look at it directly, because Jaime had gone north to stand for something, and Jaime was the father of the child turning beneath her ribs, and Jaime loved her, had always loved her, would always love her, and a woman who had buried three children and meant to keep the fourth could not afford to also believe that the man she’d loved her whole life would one day be the death of her.

She turned from the window, and put the prophecy away where she kept it, in the cold locked room at the back of her mind where she kept every truth too dangerous to look at directly, and went to see to the business of a city she had just, quite calmly, threatened to burn.